% !TeX root = report.tex
\documentclass[12pt,a4paper]{article}
\usepackage[margin=2.5cm]{geometry}
\usepackage{amsmath,amssymb,amsfonts}
\usepackage{booktabs}
\usepackage{hyperref}
\usepackage{xcolor}
\usepackage{enumitem}
\usepackage{fancyhdr}
\usepackage{mdframed}
\usepackage{listings}
\usepackage{setspace}
\usepackage{array}
\usepackage{longtable}
\usepackage{algorithm}
\usepackage{algpseudocode}
\usepackage{tikz}
\usepackage{multirow}
\usepackage{makecell}

\definecolor{blueprompt}{RGB}{0,80,160}
\definecolor{mathcolor}{RGB}{140,0,0}
\definecolor{stepcolor}{RGB}{0,120,60}
\definecolor{lightgray}{RGB}{245,245,245}
\definecolor{framecolor}{RGB}{180,180,180}
\definecolor{warningcolor}{RGB}{180,100,0}
\definecolor{novelcolor}{RGB}{80,0,120}
\definecolor{rationalecolor}{RGB}{0,100,140}
\definecolor{changecolor}{RGB}{0,100,0}

\lstset{
  backgroundcolor=\color{lightgray},
  basicstyle=\ttfamily\footnotesize,
  breaklines=true,
  frame=single,
  rulecolor=\color{framecolor}
}

\newmdenv[linecolor=blueprompt,backgroundcolor=blue!4,linewidth=1.5pt,
  innerleftmargin=10pt,innerrightmargin=10pt,
  innertopmargin=8pt,innerbottommargin=8pt]{promptbox}

\newmdenv[linecolor=stepcolor,backgroundcolor=green!4,linewidth=1.5pt,
  innerleftmargin=10pt,innerrightmargin=10pt,
  innertopmargin=8pt,innerbottommargin=8pt]{stepbox}

\newmdenv[linecolor=mathcolor,backgroundcolor=red!4,linewidth=1.5pt,
  innerleftmargin=10pt,innerrightmargin=10pt,
  innertopmargin=8pt,innerbottommargin=8pt]{mathbox}

\newmdenv[linecolor=rationalecolor,backgroundcolor=cyan!4,linewidth=1.5pt,
  innerleftmargin=10pt,innerrightmargin=10pt,
  innertopmargin=8pt,innerbottommargin=8pt]{rationalebox}

\newmdenv[linecolor=novelcolor,backgroundcolor=violet!4,linewidth=1.5pt,
  innerleftmargin=10pt,innerrightmargin=10pt,
  innertopmargin=8pt,innerbottommargin=8pt]{novelbox}

\newmdenv[linecolor=warningcolor,backgroundcolor=orange!4,linewidth=1.5pt,
  innerleftmargin=10pt,innerrightmargin=10pt,
  innertopmargin=8pt,innerbottommargin=8pt]{warningbox}

\newmdenv[linecolor=changecolor,backgroundcolor=green!3,linewidth=1.5pt,
  innerleftmargin=10pt,innerrightmargin=10pt,
  innertopmargin=8pt,innerbottommargin=8pt]{changebox}

\pagestyle{fancy}
\fancyhf{}
\rhead{\textcolor{blueprompt}{\textbf{Blueprint: Option 2 v3}}}
\lhead{\textcolor{gray}{Hard vs.\ Soft Constraint PINN --- v3 Revised}}
\rfoot{Page \thepage}

\onehalfspacing

\begin{document}

%------------------------------------------------------
\begin{titlepage}
\centering\vspace*{1.5cm}
{\LARGE\bfseries\color{blueprompt}PINN RESEARCH BLUEPRINT}\\[0.4cm]
{\large\color{gray}Research Direction 2.2 $\cdot$ Version 3 $\cdot$ June 2026}\\[1.2cm]
{\Large\bfseries
Hard vs.\ Soft Thermodynamic Constraint Enforcement 
in Physics-Informed Neural Networks\\[0.25cm]
for Sorption in Heterogeneous Porous Media}\\[1.8cm]
{\large\color{gray}
Implementation Blueprint $\cdot$ Mathematical Foundations $\cdot$\\
Unified Statistical Evaluation Harness $\cdot$ Manuscript Architecture}\\[1.2cm]

\vfill{\large\today}
\end{titlepage}

%------------------------------------------------------
\tableofcontents\newpage

%======================================================
\section{Research Vision, Gap, and Novel Contributions}
%======================================================

\subsection{The Central Scientific Question}

When physical laws are known and must be respected by a neural network
prediction, \emph{how} those laws are enforced matters. Soft constraints
penalize violations during training but allow them to persist ---
especially under extrapolation beyond the training distribution.
Hard constraints eliminate violations structurally, by construction,
but restrict the optimization landscape and may impose computational
costs. The fundamental question is:

\begin{center}
\textit{Does guaranteeing thermodynamic validity by construction
outperform penalizing violations, and at what cost in accuracy,
calibration, and computational efficiency?}
\end{center}

\subsection{Why This Has Not Been Done Before}

Hard constraint PINNs exist in fluid dynamics (topology optimization,
PDE boundary conditions) and structural mechanics. However:
\begin{itemize}
  \item No study applies hard constraints to equilibrium isotherm
    sorption models where constraints have a different mathematical
    character (algebraic inequalities and functional form requirements,
    not PDE boundary conditions).
  \item No study systematically compares \emph{each} physics constraint
    independently in a factorial design, isolating marginal contributions.
  \item No study evaluates hard vs.\ soft constraints specifically
    under \emph{extrapolation} to out-of-distribution coal ranks and
    pressures, which is the operationally critical scenario.
  \item No study combines the hard vs.\ soft comparison with calibrated
    Bayesian uncertainty quantification, asking whether hard constraints
    improve not just accuracy but predictive reliability.
\end{itemize}

\subsection{Potential Novel Contributions}

\begin{novelbox}
\textbf{Contribution 1:} Systematic hard vs.\ soft constraint
comparison for sorption isotherms in geological materials, using a
factorial ablation design that isolates and quantifies the
marginal benefit of each constraint independently and jointly.

\textbf{Contribution 2:} Three architectural strategies for
hard constraint enforcement in sorption PINNs: sigmoid-bounded
output activations (bounds), non-negative pressure pathway networks
(monotonicity), and gauge function parameterization (Sips/Arrhenius).
The Sips gauge strategy is shown to simultaneously enforce bounds,
monotonicity, and isotherm functional form as a unified architectural
choice.

\textbf{Contribution 3:} Three-tier extrapolation evaluation ---
interpolation, rank extrapolation, and pressure extrapolation ---
as a standardized benchmark for physical validity under distribution
shift, showing the greatest hard constraint advantage (to be tested!)
precisely where soft constraints fail most.

\textbf{Contribution 4:} Physical violation rate as a first-class
evaluation metric with bootstrap confidence intervals, enabling
statistical comparison of constraint enforcement strategies on
safety-relevant properties, not just accuracy.

\textbf{Contribution 5:} Constraint interaction analysis via $2\times2$
ANOVA revealing whether hard enforcement of multiple constraints
produces superadditive, additive, or antagonistic effects. Ablation
design is structured to avoid confounding of interaction terms.

\textbf{Contribution 6:} Physics parameter interpretability from the
gauge function Sips variant --- the network predicts physically
meaningful $[\hat{q}_{\max}, \hat{K}, \hat{n}]$ across coal ranks,
validated against classical isotherm literature values and SHAP/ALE
attribution. Similarly, the Arrhenius variant yields interpretable
$\Delta\hat{H}$ values validatable against calorimetric measurements.
\end{novelbox}

%======================================================
\section{Mathematical Foundations}
%======================================================

\subsection{Soft Constraint Formulation (Baseline)}

\begin{mathbox}
\textbf{Theory:} Soft constraints enter as additive penalty terms in
the loss function. They are differentiable, computationally cheap,
and allow gradient-based optimization to trade off constraint
satisfaction against data fidelity. The penalty weight $\lambda_p$
controls this trade-off but must be tuned and can never guarantee
zero violations.

\begin{align}
\mathcal{L}_{\text{total}} &= \mathcal{L}_{\text{data}}
  + \lambda_p\,\mathcal{L}_{\text{physics}} \\[4pt]
\mathcal{L}_{\text{data}} &= \mathbb{E}\!\left[
  \frac{(q - \hat{q})^2}{2\hat{\sigma}^2}
  + \tfrac{1}{2}\log\hat{\sigma}^2\right]
  \quad\text{(heteroscedastic NLL)} \\[4pt]
\mathcal{L}_{\text{physics}} &= \mathcal{L}_{\text{Sips}}
  + \mathcal{L}_{\text{bounds}}
  + \mathcal{L}_{\text{mono}}
  + \mathcal{L}_{\text{vH}} \\[4pt]
\mathcal{L}_{\text{Sips}} &= \mathrm{MSE}\!\left(\hat{q},\;
  \frac{\hat{q}_{\max}(\hat{K}P)^{\hat{n}}}{1+(\hat{K}P)^{\hat{n}}}
  \right) \\[4pt]
\mathcal{L}_{\text{bounds}} &= \mathrm{ReLU}(-\hat{q})
  + \mathrm{ReLU}(\hat{q} - \hat{q}_{\max}) \\[4pt]
\mathcal{L}_{\text{mono}} &= \frac{1}{|\mathcal{B}|}
  \sum_{i\in\mathcal{B}} \mathrm{ReLU}\!\left(
  -\frac{\partial\hat{q}}{\partial P}\bigg|_T\right) \\[4pt]
\mathcal{L}_{\text{vH}} &= \mathrm{MSE}\!\left(
  \frac{\partial\ln\hat{K}}{\partial(1/T)},\;
  -\frac{\Delta\hat{H}}{R}\right)
\end{align}
\textbf{Fundamental limitation:} $\lambda_p \to \infty$ theoretically
eliminates violations but destabilizes optimization.
In practice, violations persist at any finite $\lambda_p$.
\end{mathbox}

\subsection{Hard Constraint Strategy 1: Constrained Output Activations (Bounds)}

\begin{changebox}
\textbf{v3 correction:} The v2 formulation used $\sigma(z_s)\cdot\mathrm{softplus}(z_q)\cdot\hat{q}_{\max}$, which does not guarantee an upper bound because $\mathrm{softplus}(z_q)$ is unbounded. The formulation below is corrected: a single sigmoid output guarantees $\hat{q}\in(0,\hat{q}_{\max})$ with no additional factors needed.
\end{changebox}

\begin{mathbox}
\textbf{Theory:} Non-negativity and saturation bounds are enforced by
replacing the unbounded prediction head with a single sigmoid activation
that maps $\mathbb{R} \to (0, 1)$, scaled by a positive $\hat{q}_{\max}$.
Since $\sigma(z)\in(0,1)$ for all $z\in\mathbb{R}$, the product
$\sigma(z_q)\cdot\hat{q}_{\max}$ is strictly in $(0,\hat{q}_{\max})$
by construction, for any finite network output $z_q$.

\begin{align}
\hat{q}_{\max} &= \mathrm{softplus}\!\left(\mathrm{NN}_{q_{\max}}(\mathbf{x})\right) > 0 \\[4pt]
\hat{q} &= \hat{q}_{\max} \cdot \sigma(z_q)
  \;\in\; (0,\,\hat{q}_{\max})
\end{align}
where $z_q$ is a raw scalar network output and $\sigma$ is the sigmoid
function. This is a simpler, safer formulation than any two-factor
product: a single bounded factor is sufficient to enforce both bounds
simultaneously.

\textbf{Cost:} Negligible. One sigmoid activation per sample.
$\mathcal{L}_{\text{bounds}}$ is dropped from the loss.
\end{mathbox}

\subsection{Hard Constraint Strategy 2: Non-Negative Pressure Pathway (Monotonicity)}

\begin{changebox}
\textbf{v3 correction:} The v2 formulation used a Partially Input-Convex
Neural Network (PICNN), claiming that convexity in $P$ implies
$\partial\hat{q}/\partial P \geq 0$. This is incorrect: convexity implies
the derivative is \emph{non-decreasing}, not necessarily non-negative.
A convex function can still decrease. The formulation below replaces PICNN
with a non-negative pressure pathway architecture, which directly and
provably guarantees $\partial\hat{q}/\partial P \geq 0$.
Additionally, Swish activations on the pressure pathway are replaced by
softplus (monotone everywhere); Swish has a non-monotone region near
$x \approx -1.3$ that can break the monotonicity guarantee.
\end{changebox}

\begin{mathbox}
\textbf{Theory:} A sufficient condition for $\partial\hat{q}/\partial P \geq 0$
is that every pathway through which $P$ contributes to $\hat{q}$
uses (a) non-negative weights and (b) monotone increasing activations.
Under these two conditions, increasing $P$ cannot decrease $\hat{q}$,
regardless of the network depth.

For a single pressure-pathway layer:
\begin{align}
\hat{q}(P, \mathbf{x}_c) &= a(\mathbf{x}_c)
  + \mathrm{softplus}(w \cdot P + g(\mathbf{x}_c)) \\[4pt]
\frac{\partial\hat{q}}{\partial P} &= w \cdot \sigma(wP + g(\mathbf{x}_c)) \geq 0
  \quad\text{iff } w \geq 0
\end{align}
Since $\sigma(\cdot) > 0$ always, the sign of $\partial\hat{q}/\partial P$
is determined entirely by the sign of $w$. Non-negativity is enforced via
reparameterization during every forward pass:
\begin{align}
w &= \mathrm{softplus}(w_{\text{raw}})
  \quad (w_{\text{raw}} \in \mathbb{R},\; w > 0)
\end{align}

For a multi-layer pressure pathway with $L$ layers:
\begin{align}
z_P^{(0)} &= P \\[2pt]
z_P^{(\ell+1)} &= \mathrm{softplus}\!\left(
  \underbrace{W_P^{(\ell)}}_{\geq 0} z_P^{(\ell)}
  + W_c^{(\ell)} \mathbf{x}_c + b^{(\ell)}\right) \\[4pt]
\hat{q} &= W_{\text{out}} z_P^{(L)} + a(\mathbf{x}_c)
\end{align}
where $W_P^{(\ell)} = \mathrm{softplus}(W_{P,\text{raw}}^{(\ell)}) \geq 0$
is recomputed at every forward pass.

\textbf{Critical implementation requirement:} Pressure $P$ must enter the
network \emph{exclusively} through this constrained pathway. If $P$ is
also passed through any standard (unconstrained) layer, the guarantee is
lost. The compositional features $\mathbf{x}_c$ may enter freely through
unconstrained pathways.

\textbf{Activation choice:} softplus or sigmoid only on the pressure
pathway. Swish is not permitted on this pathway because it is non-monotone
for $x \lesssim -1.3$, which would break the monotonicity guarantee.

\textbf{Cost:} Moderate. Softplus reparameterization added per layer;
one additional backward pass for Jacobian verification.
$\mathcal{L}_{\text{mono}}$ is dropped.
\end{mathbox}

\subsection{Hard Constraint Strategy 3: Gauge Function Parameterization (Sips)}

\begin{changebox}
\textbf{v3 note:} The Sips gauge variant simultaneously enforces bounds,
monotonicity, and Sips functional form, provided $\hat{q}_{\max}, \hat{K},
\hat{n} > 0$ and raw physical pressure $P > 0$ is used. This makes separate
hard-bounds (Strategy 1) and hard-monotonicity (Strategy 2) components
redundant within this variant. The ablation design accounts for this explicitly.
\end{changebox}

\begin{mathbox}
\textbf{Theory:} The Sips isotherm shape is enforced by structuring
the network to predict isotherm \emph{parameters} from coal composition,
then computing the isotherm analytically. The functional form becomes
the output layer itself.

\begin{align}
\hat{q}(\mathbf{x}_c, P) &= \frac{\hat{q}_{\max}(\mathbf{x}_c)
  \cdot [\hat{K}(\mathbf{x}_c)\cdot P]^{\hat{n}(\mathbf{x}_c)}}
  {1 + [\hat{K}(\mathbf{x}_c)\cdot P]^{\hat{n}(\mathbf{x}_c)}} \\[6pt]
\hat{q}_{\max} &= \mathrm{softplus}(\mathrm{NN}_1(\mathbf{x}_c)) > 0 \\[4pt]
\hat{K} &= \mathrm{softplus}(\mathrm{NN}_2(\mathbf{x}_c)) > 0 \\[4pt]
\hat{n} &= \mathrm{softplus}(\mathrm{NN}_3(\mathbf{x}_c)) > 0
\end{align}

\textbf{Structural guarantees (given $P > 0$):}
\begin{itemize}[noitemsep]
  \item $\hat{q} > 0$ (positivity): numerator always positive
  \item $\hat{q} < \hat{q}_{\max}$ (saturation bound): Sips ratio $< 1$
  \item $\partial\hat{q}/\partial P \geq 0$ (monotonicity): Sips is
    monotone increasing in $P$ for all positive parameters
\end{itemize}
The Sips gauge variant therefore subsumes hard-bounds and hard-monotonicity
as a special case. Combining it with separate bounds or monotonicity
enforcement (V7 in v2 design) is architecturally redundant.

\textbf{Implementation requirement:} Raw physical pressure $P > 0$ (in MPa)
must be used in the Sips layer, not a scaled or log-transformed version.
Assert $P > 0$ in the data pipeline. Pressure enters \emph{only} at the
output Sips layer, not into the encoder.

\textbf{Note:} The predicted parameters $[\hat{q}_{\max}, \hat{K}, \hat{n}]$
are directly physically interpretable and validatable against classical
isotherm literature. $\mathcal{L}_{\text{Sips}}$ is dropped.
\end{mathbox}

\subsection{Hard Constraint Strategy 4: Arrhenius Reparameterization (Van't~Hoff)}

\begin{mathbox}
\textbf{Theory:} Van't~Hoff consistency requires $K(T)$ to follow
Arrhenius kinetics. Rather than penalizing deviations, the network is
parameterized to produce $K(T)$ that satisfies it exactly by definition.

\begin{align}
\hat{K}(T, \mathbf{x}_c) &= \hat{K}_0(\mathbf{x}_c)
  \cdot \exp\!\left(-\frac{\Delta\hat{H}(\mathbf{x}_c)}{RT}\right)
\end{align}
where $\hat{K}_0(\mathbf{x}_c) = \mathrm{softplus}(\mathrm{NN}_{K_0}(\mathbf{x}_c)) > 0$
and $\Delta\hat{H}(\mathbf{x}_c) = -\mathrm{softplus}(\mathrm{NN}_H(\mathbf{x}_c)) < 0$
(exothermic adsorption, physically required for physisorption).

By construction, $\partial\ln\hat{K}/\partial(1/T) = -\Delta\hat{H}/R$
exactly. The predicted $\Delta\hat{H}$ values (in kJ/mol) are physically
interpretable adsorption enthalpies, validatable against calorimetric
measurements ($-15$ to $-35$\,kJ/mol for methane-coal).
$\mathcal{L}_{\text{vH}}$ is dropped.
\end{mathbox}

\subsection{Constraint Subsumption Map}

\begin{rationalebox}
A key structural insight emerging from the corrected formulations is that
the Sips gauge architecture (Strategy 3) subsumes the bounds guarantee
(Strategy 1) and the pressure-monotonicity guarantee (Strategy 2),
as long as all Sips parameters are positive and raw $P > 0$ is used.
This has direct implications for which ablation variants are independent.

\begin{center}
\small
\begin{tabular}{lcccc}
\toprule
\textbf{Strategy} & \textbf{Bounds} & \textbf{Mono} & \textbf{Sips form} & \textbf{Van't~Hoff} \\
\midrule
Strategy 1 (sigmoid head) & \checkmark & --- & --- & --- \\
Strategy 2 (NN$^+$ pathway) & --- & \checkmark & --- & --- \\
Strategy 3 (Sips gauge) & \checkmark & \checkmark & \checkmark & --- \\
Strategy 4 (Arrhenius) & --- & --- & --- & \checkmark \\
\bottomrule
\end{tabular}
\end{center}

Strategies 3 and 4 together (Sips gauge + Arrhenius) enforce all four
physical constraints and form the all-hard variant~V2.
\end{rationalebox}

\subsection{Revised Factorial Ablation Design}

\begin{changebox}
\textbf{v3 revision:} The ablation design is revised to correct two issues
from v2: (a) the old V8 (Hard-\{S+B+M\}) was redundant because the Sips
gauge already enforces bounds and monotonicity; (b) the Sips$\times$vH
interaction was confounded in the old V2/V8 because bounds and monotonicity
were simultaneously hard. The new V8 (Hard-\{S+vH\}) enables a clean
Sips$\times$vH interaction test using \{V1, V5, V6, V8\}, analogous to
the clean Bounds$\times$Monotonicity test using \{V1, V3, V4, V7\}.
\end{changebox}

\begin{mathbox}
\textbf{Design:} Eight scientifically motivated variants preserving all
main effects and the two primary interaction subspaces:

\begin{center}
\small
\begin{tabular}{llccccl}
\toprule
\textbf{ID} & \textbf{Variant} & \textbf{Bounds} & \textbf{Mono} &
  \textbf{Sips} & \textbf{vH} & \textbf{Purpose} \\
\midrule
V1 & All-soft      & S & S & S & S & Baseline \\
V2 & All-hard      & H$^*$ & H$^*$ & H & H & Maximum enforcement \\
V3 & Hard-B        & H & S & S & S & Isolate bounds \\
V4 & Hard-M        & S & H & S & S & Isolate monotonicity \\
V5 & Hard-S        & H$^*$ & H$^*$ & H & S & Isolate Sips (subsumes B+M) \\
V6 & Hard-vH       & S & S & S & H & Isolate thermal \\
V7 & Hard-\{B+M\}  & H & H & S & S & Joint bounds+mono; B$\times$M interaction \\
V8 & Hard-\{S+vH\} & H$^*$ & H$^*$ & H & H & Sips+thermal; S$\times$vH interaction \\
\bottomrule
\multicolumn{7}{l}{\small $^*$ Enforced implicitly by the Sips gauge architecture.}\\
\multicolumn{7}{l}{\small H = hard, S = soft.}\\
\end{tabular}
\end{center}

\textbf{Interaction tests enabled:}
\begin{itemize}[noitemsep]
  \item Bounds $\times$ Monotonicity: \{V1, V3, V4, V7\} --- clean $2\times2$
  \item Sips $\times$ Van't~Hoff: \{V1, V5, V6, V8\} --- clean $2\times2$
\end{itemize}

\textbf{Note on V5:} Because the Sips gauge subsumes bounds and
monotonicity, V5 is not a pure single-constraint isolate. This should
be explicitly acknowledged in the manuscript. V5 tests the effect of
the Sips gauge architecture as a whole (which includes implicit B+M
enforcement), not the Sips functional form constraint in isolation.
This is an inherent consequence of the architectural approach, not a
design flaw, and is scientifically informative in its own right.
\end{mathbox}

%======================================================
\section{Implementation Pipeline}
%======================================================

\begin{rationalebox}
\textbf{Pipeline design principle:} The evaluation harness is built
\emph{before} any model, not after. Every model variant plugs into
the same harness automatically. This prevents retrofitting, ensures
fair comparison, and eliminates the risk of adapting evaluation
procedures to favour certain results.
\end{rationalebox}

\subsection*{Pipeline Overview}

\begin{center}
\small
\begin{tabular}{clll}
\toprule
\textbf{Step} & \textbf{Component} & \textbf{Produces} & \textbf{Used by} \\
\midrule
1 & Data pipeline + 3-tier splits & Loaders, scalers & All steps \\
2 & Evaluation harness & Stats framework & Steps 4--10 \\
3 & Soft baseline + $\lambda_p$ sweep & V1 + sensitivity & Manuscript \S3.1 \\
4 & Hard bounds (V3) & Architecture A & Harness \\
5 & Hard monotonicity (V4) & Architecture B & Harness \\
6 & Hard Sips --- gauge (V5) & Architecture C & Harness \\
7 & Hard vH --- Arrhenius (V6) & Architecture D & Harness \\
8 & Joint variants V2, V7, V8 & Final 3 variants & Harness \\
9 & Full ablation + interaction & All 8 variants & Manuscript \S3.2--3.5 \\
10 & UQ calibration & Calibrated models & Manuscript \S3.6 \\
11 & Interpretability & SHAP + param plots & Manuscript \S3.7 \\
\bottomrule
\end{tabular}
\end{center}

%======================================================
\section{Step-by-Step Implementation}
%======================================================

\subsection{Step 1: Data Pipeline with Three-Tier Evaluation}

\begin{rationalebox}
\textbf{Rationale:} Extrapolation performance under distribution shift
is the primary scientific argument for hard constraints. Three evaluation
tiers must be defined from the start so every model is evaluated
identically. GroupKFold at the experiment level prevents sample-level
leakage across all splits.

\textbf{v3 addition:} An explicit assertion that $P > 0$ for all samples
is required, since the Sips gauge monotonicity guarantee depends on raw
positive pressure. Pressure values at or below zero are physically
meaningless in this context and would produce NaN gradients in
$(\hat{K} P)^{\hat{n}}$.
\end{rationalebox}

\begin{stepbox}
\textbf{What it achieves:}
\begin{itemize}[noitemsep]
  \item Experiment-level GroupShuffleSplit (no leakage)
  \item Log target transform + robust scaling (fitted on train only)
  \item 12-feature engineering (same as original study)
  \item 5-fold GroupKFold for cross-validation
  \item Three evaluation tiers: T1 (interpolation), T2 (rank extrap),
    T3 (pressure extrap)
  \item Raw pressure tensor preserved (unscaled MPa) for Sips output layer
  \item Assert $P > 0$ for all samples before training
\end{itemize}
\end{stepbox}

\begin{promptbox}
\textbf{Requirements --- Step 1:}

\texttt{Implement a complete data pipeline in PyTorch for methane coal
sorption modeling.}

\texttt{INPUT CSV columns: experiment\_id, coal\_rank (anthracite/lv-bitu/
mv-bitu/hv-bitu/lignite), temperature\_K, pressure\_MPa, moisture\_wt,
ash\_wt, volatile\_wt, CH4\_adsorption\_m3t}

\texttt{SPLITS (implement all, no overlap):}
\begin{itemize}[noitemsep]
  \item \texttt{Main split: GroupShuffleSplit 80/20 by experiment\_id,
    stratify by joint temperature-pressure bin within groups}
  \item \texttt{T1 (interpolation): standard 20\% held-out test set}
  \item \texttt{T2 (rank extrapolation): ALL lignite experiments
    completely withheld from training. If lignite in test set, move
    to T2. T2 evaluated separately.}
  \item \texttt{T3 (pressure extrapolation): from T1, isolate samples
    with P > 8.0 MPa (top 10th percentile). Evaluate models on these
    without retraining.}
  \item \texttt{5-fold CV: GroupKFold(n\_splits=5) on training set,
    group=experiment\_id}
\end{itemize}

\texttt{FEATURES: engineer Tr=T/190.6, Pr=P/4.60, FC=100-M-A-VM,
OM=(VM+FC)/100, beta=1/(8.314*T), PT=P*T, MV=moisture*volatile.
Total: 12 features. Robust scale (median/IQR) fitted on train fold only.}

\texttt{PRESSURE HANDLING (critical for Sips gauge):}
\begin{itemize}[noitemsep]
  \item \texttt{Return raw\_pressure tensor (unscaled MPa) alongside
    scaled feature tensor. The Sips output layer uses raw pressure.}
  \item \texttt{assert (raw\_pressure > 0).all(), "Sips monotonicity
    guarantee requires P > 0 for all samples"}
  \item \texttt{Log minimum pressure in dataset; warn if any P < 0.01 MPa}
\end{itemize}

\texttt{RETURN: DataLoaders (train/val/T1/T2/T3), scaler object,
feature\_names list, fold\_indices for CV, raw\_pressure tensor
(unscaled MPa), experiment\_id tensor.}

\texttt{VERIFICATION: assert zero experiment overlap across splits.
Assert T2 contains only lignite. Assert T3 subset of T1. Assert P > 0.
Print split sizes and rank distribution per split.}
\end{promptbox}

\subsection{Step 2: Unified Evaluation Harness}

\begin{rationalebox}
\textbf{Rationale:} Building the evaluation harness before any model
ensures all variants are compared on identical metrics computed
identically. Statistical rigor (bootstrap CIs, corrected tests) is
applied from the first result. Physical violation rates with uncertainty
bounds are treated as first-class results.
\end{rationalebox}

\begin{promptbox}
\textbf{Requirements --- Step 2:}

\texttt{Implement a unified evaluation harness for PINN constraint
comparison.}

\texttt{FUNCTION: evaluate\_model(model, data\_dict, n\_bootstrap=100,
n\_mc=100, alpha=0.05)}

\texttt{ACCURACY METRICS (compute on T1, T2, T3 separately):}
\begin{itemize}[noitemsep]
  \item \texttt{R2, RMSE, MAE, MaxAE on original scale (back-transform)}
  \item \texttt{Median relative error (\%), mean relative error (\%)}
  \item \texttt{95\% bootstrap CI for each metric (stratified by coal rank)}
\end{itemize}

\texttt{PHYSICAL VIOLATION METRICS (primary):}
\begin{itemize}[noitemsep]
  \item \texttt{neg\_rate: fraction samples with q\_hat < 0 (+ 95\% CI)}
  \item \texttt{sat\_rate: fraction with q\_hat > q\_max (+ 95\% CI)}
  \item \texttt{nonmono\_rate: fraction with dq/dP < 0 at any P in
    linspace(0.01, 12, 50) per sample (+ 95\% CI)}
  \item \texttt{vH\_mae: mean |d(lnK)/d(1/T) + dH/R| (+ 95\% CI)}
  \item \texttt{For hard constraints: assert 0.0 and raise ValueError
    if non-zero (structural guarantee must hold)}
\end{itemize}

\texttt{UQ METRICS (MC Dropout, 100 passes):}
\begin{itemize}[noitemsep]
  \item \texttt{sigma\_epi, sigma\_ale, sigma\_total}
  \item \texttt{ECE (4 confidence levels: 68, 90, 95, 99\%)}
  \item \texttt{rho\_s: Spearman(|error|, sigma\_total)}
  \item \texttt{NLL, CRPS, Sharpness (95\% PI width)}
  \item \texttt{Coverage at 68\%, 90\%, 95\%, 99\%}
\end{itemize}

\texttt{COMPUTATIONAL METRICS:}
\begin{itemize}[noitemsep]
  \item \texttt{train\_time\_sec, inference\_time\_ms\_per\_sample}
  \item \texttt{n\_parameters, memory\_MB (peak GPU)}
  \item \texttt{epochs\_to\_R2\_085, epochs\_to\_R2\_090}
\end{itemize}

\texttt{STATISTICAL COMPARISON FUNCTIONS:}
\begin{itemize}[noitemsep]
  \item \texttt{compare\_variants(v1, v2, metric): bootstrap paired
    t-test + Cohen's d + p-value}
  \item \texttt{bonferroni\_correct(p\_values, n=28): corrected alpha
    = 0.05/28 = 0.0018}
  \item \texttt{anova\_2x2(results, factor1, factor2, metric)}
  \item \texttt{mcnemar\_test(violations\_v1, violations\_v2)}
\end{itemize}

\texttt{OUTPUT: dict with all metrics + CIs, auto-generate LaTeX table.
Log everything to CSV for reproducibility.}
\end{promptbox}

\subsection{Step 3: Soft Baseline + $\lambda_p$ Sensitivity Analysis}

\begin{rationalebox}
\textbf{Rationale:} Before declaring hard constraints superior, we must
rule out that soft baseline underperforms simply due to poor $\lambda_p$
tuning. The sensitivity sweep is mandatory to make the hard vs.\ soft
comparison credible.
\end{rationalebox}

\begin{promptbox}
\textbf{Requirements --- Step 3:}

\texttt{Implement the soft-constraint PINN baseline (V1) and lambda\_p
sensitivity sweep.}

\texttt{ARCHITECTURE (V1 baseline):}
\begin{itemize}[noitemsep]
  \item \texttt{Input: 12 features (including scaled pressure)}
  \item \texttt{5 residual blocks [512,1024,512,256,128],
    LayerNorm (NOT BatchNorm), Swish activations, dropout=0.1}
  \item \texttt{3 output heads: mean (softplus for positivity --- soft),
    log-variance (aleatoric), physics params [q\_max,K,n] via softplus}
  \item \texttt{All 4 physics losses active as soft penalties}
\end{itemize}

\texttt{LAMBDA\_P SENSITIVITY SWEEP:}
\begin{itemize}[noitemsep]
  \item \texttt{Train V1 with lambda\_p in [0.0, 0.01, 0.05, 0.1,
    0.2, 0.5, 1.0, 2.0]}
  \item \texttt{For each: 3 seeds, 5-fold CV}
  \item \texttt{Record: R2, RMSE, violation rates, ECE}
  \item \texttt{Select lambda\_p* = argmin validation RMSE}
  \item \texttt{Plot: dual-axis RMSE vs lambda\_p and violation rate
    vs lambda\_p (Pareto trade-off)}
  \item \texttt{KEY RESULT: is there a lambda\_p achieving zero
    violations AND high R2? (Expected: no)}
\end{itemize}

\texttt{TRAINING: AdamW, lr=1e-3, wd=1e-5, early stopping patience=100,
max\_epochs=1000. Log per epoch: loss, R2\_val, violation\_rates.}

\texttt{Plug into harness. Store in results\_dict['V1'].}
\end{promptbox}

\subsection{Step 4: Hard Bounds Variant (V3)}

\begin{rationalebox}
\textbf{Rationale:} Physical bounds are the simplest constraint to
enforce exactly and have the smallest computational overhead. Even this
trivial hard constraint establishing measurable benefit over soft
enforcement anchors the argument for the full study.
\end{rationalebox}

\begin{stepbox}
\textbf{What it achieves:} V3 (Hard-B) with corrected sigmoid-only output
head that structurally maps to $(0, \hat{q}_{\max})$.
$\mathcal{L}_{\text{bounds}}$ removed; all other losses unchanged.
\end{stepbox}

\begin{promptbox}
\textbf{Requirements --- Step 4:}

\texttt{Implement V3 (Hard-B): soft baseline with hard physical bounds.
Modify ONLY the mean prediction head from V1:}

\texttt{class HardBoundsHead(nn.Module):}
\begin{itemize}[noitemsep]
  \item \texttt{head\_qmax: Linear -> softplus -> q\_max > 0}
  \item \texttt{head\_q: Linear -> raw scalar z\_q}
  \item \texttt{forward(h):}
  \begin{itemize}[noitemsep]
    \item \texttt{q\_max = softplus(head\_qmax(h))}
    \item \texttt{z\_q = head\_q(h)}
    \item \texttt{q\_hat = q\_max * torch.sigmoid(z\_q)}
    \item \texttt{\# q\_hat is strictly in (0, q\_max) for all finite z\_q}
    \item \texttt{return q\_hat, q\_max}
  \end{itemize}
\end{itemize}

\texttt{VERIFICATION:}
\begin{itemize}[noitemsep]
  \item \texttt{verify\_bounds(model, loader): assert (q\_hat > 0).all()
    and (q\_hat < q\_max).all()}
  \item \texttt{Run on 1000 random inputs including extreme P values}
  \item \texttt{Raise ValueError with violation details if assertion fails}
\end{itemize}

\texttt{Remove L\_bounds. Keep L\_Sips + L\_mono + L\_vH soft.
Plug into harness. Store in results\_dict['V3'].}
\end{promptbox}

\subsection{Step 5: Hard Monotonicity via Non-Negative Pressure Pathway (V4)}

\begin{rationalebox}
\textbf{Rationale:} Monotonic pressure-adsorption is a thermodynamic
law (Le~Chatelier). Violations are not just statistically undesirable
--- they are physically impossible. Soft constraints can produce
non-monotonic isotherms for rare compositional inputs at high pressure,
exactly the extrapolation scenario most relevant for reservoir assessment.

The corrected architecture directly guarantees $\partial\hat{q}/\partial P \geq 0$
through non-negative weights and monotone activations on the pressure
pathway, without relying on convexity arguments.
\end{rationalebox}

\begin{promptbox}
\textbf{Requirements --- Step 5:}

\texttt{Implement V4 (Hard-M): non-negative pressure pathway for hard
monotonicity enforcement.}

\texttt{ARCHITECTURE:}
\begin{itemize}[noitemsep]
  \item \texttt{Compositional encoder: standard 5 residual blocks
    [512,1024,512,256,128] with Swish. Input: 11 compositional features
    (no pressure). Output: h\_comp (128-dim)}
  \item \texttt{Pressure pathway: separate 3-layer network with
    hidden dims [64, 32, 16]. CRITICAL:}
  \begin{itemize}[noitemsep]
    \item \texttt{All weight matrices on pressure-input connections
      reparameterized: W = F.softplus(W\_raw) >= 0}
    \item \texttt{Activations: softplus ONLY (not Swish, not ReLU).
      Swish is non-monotone and breaks the guarantee.}
    \item \texttt{Pressure P enters ONLY this pathway, nowhere else.}
  \end{itemize}
  \item \texttt{Output: q\_hat = linear(concat(h\_comp, h\_pressure))}
\end{itemize}

\texttt{REPARAMETERIZATION (critical --- must be in forward pass):}
\begin{itemize}[noitemsep]
  \item \texttt{def forward(self, x\_comp, P\_raw):}
  \item \texttt{W\_eff = F.softplus(self.W\_raw) \# recompute each pass}
  \item \texttt{h\_P = softplus(W\_eff @ P\_raw + self.b)}
  \item \texttt{Do NOT store softplus(W\_raw) as a parameter.}
\end{itemize}

\texttt{VERIFICATION:}
\begin{itemize}[noitemsep]
  \item \texttt{verify\_monotonicity(model, n\_comp=500, n\_P=50):}
  \item \texttt{Sample 500 random coal compositions from training dist}
  \item \texttt{Sweep P in linspace(0.01, 15, 50)}
  \item \texttt{dq\_dP = torch.autograd.grad(q\_hat.sum(), P\_raw)[0]}
  \item \texttt{assert (dq\_dP >= -1e-6).all()}
  \item \texttt{Report: max violation magnitude, fraction negative}
\end{itemize}

\texttt{Remove L\_mono. Keep L\_Sips + L\_bounds + L\_vH soft.
Plug into harness. Store in results\_dict['V4'].}
\end{promptbox}

\subsection{Step 6: Gauge Function Sips Architecture (V5)}

\begin{rationalebox}
\textbf{Rationale:} The Sips gauge variant enforces the functional form
of the pressure-adsorption relationship and, as a structural consequence,
also guarantees bounds and monotonicity. The encoder predicts physically
interpretable parameters $[\hat{q}_{\max}, \hat{K}, \hat{n}]$ from
composition alone, enabling external validation against classical isotherm
literature and SHAP attribution per parameter.

The manuscript should explicitly acknowledge that V5 is not a single-constraint
isolate but a composite architectural choice. The scientific value of V5
is both in testing the Sips functional form constraint and in demonstrating
the interpretability bonus of the gauge approach.
\end{rationalebox}

\begin{promptbox}
\textbf{Requirements --- Step 6:}

\texttt{Implement V5 (Hard-S): gauge function Sips PINN.}

\texttt{class SipsGaugePINN(nn.Module):}
\begin{itemize}[noitemsep]
  \item \texttt{Shared encoder: 5 residual blocks [512,1024,512,256,128],
    input=11 compositional features (NO pressure in encoder)}
  \item \texttt{head\_qmax: Linear(128,1) -> softplus -> q\_max > 0}
  \item \texttt{head\_K: Linear(128,1) -> softplus -> K > 0}
  \item \texttt{head\_n: Linear(128,1) -> softplus -> n > 0}
  \item \texttt{head\_logvar: Linear(128,1) -> log\_sigma2}
  \item \texttt{forward(x\_comp, P\_raw):}
  \begin{itemize}[noitemsep]
    \item \texttt{assert (P\_raw > 0).all(), "Sips requires P > 0"}
    \item \texttt{h = encoder(x\_comp)}
    \item \texttt{q\_max, K, n = heads(h)}
    \item \texttt{KP\_n = (K * P\_raw).pow(n)}
    \item \texttt{q\_hat = q\_max * KP\_n / (1 + KP\_n)}
    \item \texttt{return q\_hat, log\_sigma2, q\_max, K, n}
  \end{itemize}
\end{itemize}

\texttt{STRUCTURAL VERIFICATION:}
\begin{itemize}[noitemsep]
  \item \texttt{Assert q\_hat > 0 and q\_hat < q\_max for all samples}
  \item \texttt{Assert dq/dP >= 0 for linspace(0.01, 15, 50)}
  \item \texttt{These should hold by construction; any violation
    indicates a bug (not a model limitation)}
\end{itemize}

\texttt{PARAMETER ANALYSIS:}
\begin{itemize}[noitemsep]
  \item \texttt{Compute mean q\_max, K, n per coal rank}
  \item \texttt{Assert q\_max in [10,40] m3/t, K in [0.1,5] MPa-1,
    n in [0.3,2.0]}
  \item \texttt{SHAP on each parameter head separately}
\end{itemize}

\texttt{Remove L\_Sips, L\_bounds, L\_mono (all subsumed).
Keep L\_vH soft. Plug into harness. Store in results\_dict['V5'].}
\end{promptbox}

\subsection{Step 7: Arrhenius Reparameterization (V6)}

\begin{rationalebox}
\textbf{Rationale:} Van't~Hoff consistency governs temperature dependence
--- critical for predictions at temperatures outside the narrow 20--50$^\circ$C
training range. The Arrhenius reparameterization produces interpretable
$\Delta\hat{H}$ values validatable against calorimetric literature.
\end{rationalebox}

\begin{promptbox}
\textbf{Requirements --- Step 7:}

\texttt{Implement V6 (Hard-vH): Arrhenius K(T) reparameterization.
Modify the physics parameter head from V1:}

\texttt{class ArrheniusHead(nn.Module):}
\begin{itemize}[noitemsep]
  \item \texttt{head\_K0: Linear(128,1) -> softplus -> K0 > 0}
  \item \texttt{head\_dH: Linear(128,1) -> -softplus(x) -> dH < 0
    (exothermic)}
  \item \texttt{forward(h, T\_raw):}
    \texttt{K = K0 * torch.exp(dH / (8.314 * T\_raw))}
    \texttt{(note: dH < 0, so exponent is negative as required)}
  \item \texttt{Return K, K0, dH (convert to kJ/mol for display)}
\end{itemize}

\texttt{VERIFICATION:}
\begin{itemize}[noitemsep]
  \item \texttt{Compute d(lnK)/d(1/T) numerically for 100 samples}
  \item \texttt{Assert max|d(lnK)/d(1/T) + dH/R| < 1e-4}
  \item \texttt{Report dH distribution per coal rank in kJ/mol}
  \item \texttt{Assert 95\% of dH in [-60, -5] kJ/mol}
\end{itemize}

\texttt{Compare dH to Busch \& Gensterblum (2011):
anthracite ~-22, bituminous ~-18, lignite ~-14 kJ/mol.
Plug into harness. Store in results\_dict['V6'].}
\end{promptbox}

\subsection{Step 8: Joint Variants V2, V7, V8}

\begin{rationalebox}
\textbf{Rationale:} Joint variants test whether hard constraints combine
additively or otherwise. V7 (Hard-\{B+M\}) is the most practically
relevant: zero-violation guarantee on the two most commonly violated
constraints at minimal cost. V8 (Hard-\{S+vH\}) enables the clean
Sips$\times$vH interaction test. V2 (all-hard) tests whether maximum
enforcement hurts accuracy.
\end{rationalebox}

\begin{promptbox}
\textbf{Requirements --- Step 8:}

\texttt{V7 (Hard-B+M): combine HardBoundsHead (Step 4) with non-negative
pressure pathway (Step 5). Drop L\_bounds and L\_mono.}

\texttt{V8 (Hard-S+vH): SipsGaugePINN (Step 6) with ArrheniusHead
(Step 7). All four physical losses dropped. L\_data only.
Note: Sips gauge already enforces bounds and mono implicitly.}

\texttt{V2 (All-hard): identical to V8. Sips + Arrhenius together
enforce all four constraints. L\_data only.}

\texttt{For each variant: run all structural verification functions.
If V2 shows degraded accuracy vs V1, report this honestly.
Store in results\_dict.}

\texttt{NOTE: V2 and V8 are equivalent under this ablation design.
Run as one variant. Report in the manuscript that the all-hard model
is achieved by combining Sips gauge + Arrhenius, which is also V8.}
\end{promptbox}

\subsection{Step 9: Full Ablation + Interaction Analysis}

\begin{promptbox}
\textbf{Requirements --- Step 9:}

\texttt{MAIN EFFECTS TABLE: for all 8 variants on T1, T2, T3:}
\begin{itemize}[noitemsep]
  \item \texttt{R2 (mean +/- 95\% CI), RMSE, MAE, MaxAE}
  \item \texttt{violation rates with CIs (hard variants must show 0.0)}
  \item \texttt{ECE, rho\_s, NLL, train time, parameters, memory}
\end{itemize}

\texttt{PAIRWISE STATISTICAL TESTS (28 pairs):}
\begin{itemize}[noitemsep]
  \item \texttt{Bootstrap paired t-test on RMSE (100 bootstrap samples)}
  \item \texttt{Bonferroni correction: alpha = 0.05/28 = 0.0018}
  \item \texttt{Cohen's d for effect size}
  \item \texttt{McNemar's test for violation rate proportions}
  \item \texttt{Report as significance matrix (heatmap)}
\end{itemize}

\texttt{ANOVA INTERACTION ANALYSIS:}
\begin{itemize}[noitemsep]
  \item \texttt{Bounds x Mono: use V1(S,S), V3(H,S), V4(S,H), V7(H,H)}
  \item \texttt{Sips x vH: use V1(S,S), V5(H,S), V6(S,H), V8(H,H)}
  \item \texttt{Interaction delta: (V7-V1) - (V3-V1) - (V4-V1)
    and (V8-V1) - (V5-V1) - (V6-V1)}
  \item \texttt{delta > 0: superadditive; delta < 0: antagonistic}
\end{itemize}

\texttt{EXTRAPOLATION ADVANTAGE:}
\begin{itemize}[noitemsep]
  \item \texttt{R2 retention ratio: R2\_T2/R2\_T1 and R2\_T3/R2\_T1
    per variant}
  \item \texttt{Plot: retention ratio vs violation rate (expected
    negative correlation)}
\end{itemize}
\end{promptbox}

\subsection{Step 10: Calibrated UQ Across All Variants}

\begin{rationalebox}
\textbf{Rationale:} A key hypothesis is that hard constraints, by narrowing
the admissible solution space, produce lower epistemic uncertainty. If
confirmed, this provides a mechanistic explanation for ensemble collapse
observed in the original paper, connecting the two studies theoretically.
\end{rationalebox}

\begin{promptbox}
\textbf{Requirements --- Step 10:}

\texttt{For each variant (100 MC Dropout passes):}
\begin{itemize}[noitemsep]
  \item \texttt{Compute epistemic\_fraction\% per variant}
  \item \texttt{Fit tau via ECE minimization on validation set}
  \item \texttt{Evaluate calibrated ECE, rho\_s, NLL, CRPS on T1, T2, T3}
\end{itemize}

\texttt{KEY PLOT: epistemic\_fraction vs constraint\_hardness\_score
(score = number of hard constraints active, 0--4).
Expected: monotone decreasing relationship.}

\texttt{CALIBRATION UNDER EXTRAPOLATION:}
\begin{itemize}[noitemsep]
  \item \texttt{ECE on T1 and T2 separately per variant}
  \item \texttt{Report ECE\_T2/ECE\_T1 ratio (hard constraints expected
    to maintain better calibration under shift)}
\end{itemize}
\end{promptbox}

\subsection{Step 11: Interpretability and Physical Parameter Validation}

\begin{rationalebox}
\textbf{Rationale:} The gauge function variant (V5) produces parameters
validatable against external literature --- impossible for standard PINNs.
SHAP per parameter head reveals what compositional features control
maximum capacity, equilibrium constant, and heterogeneity. This is new
scientific knowledge about coal sorption, not just model evaluation.
\end{rationalebox}

\begin{promptbox}
\textbf{Requirements --- Step 11:}

\texttt{GAUGE FUNCTION PARAMETER ANALYSIS (V5):}
\begin{itemize}[noitemsep]
  \item \texttt{q\_max, K, n distributions per coal rank (violin plots)}
  \item \texttt{Assert Pearson corr(q\_max, volatile) < -0.3}
  \item \texttt{SHAP on head\_qmax, head\_K, head\_n separately}
  \item \texttt{Compare to classical Langmuir fits from Step 3}
\end{itemize}

\texttt{ARRHENIUS dH VALIDATION (V6):}
\begin{itemize}[noitemsep]
  \item \texttt{dH mean +/- std per coal rank in kJ/mol}
  \item \texttt{Compare to Busch \& Gensterblum (2011)}
  \item \texttt{One-way ANOVA + Tukey HSD across ranks}
\end{itemize}

\texttt{GLOBAL SHAP COMPARISON (V1 vs V5):}
\begin{itemize}[noitemsep]
  \item \texttt{Mean |SHAP| feature importance ranking for both}
  \item \texttt{Rank correlation between V1 and V5 importance}
\end{itemize}
\end{promptbox}

%======================================================
\section{Key Connections Between Steps}
%======================================================

\begin{center}
\small
\begin{tabular}{p{3cm}p{4cm}p{4.5cm}}
\toprule
\textbf{Step produces} & \textbf{Used in} & \textbf{How} \\
\midrule
Step 1: T1/T2/T3 splits & All steps & Identical evaluation on all tiers \\
Step 1: raw\_pressure & Steps 6,7,8 & Required by Sips gauge output layer \\
Step 2: Harness & Steps 3--11 & evaluate\_model() for every variant \\
Step 3: $\lambda_p^*$ & Steps 4--8 & All soft components use optimal $\lambda_p^*$ \\
Step 3: V1 results & Step 9 & Baseline for all pairwise comparisons \\
Step 5: NN$^+$ pathway & Step 8 (V7) & Reused in joint variant \\
Step 6: Gauge code & Step 8 (V2/V8) & Reused in joint variants \\
Step 7: Arrhenius & Step 8 (V2/V8) & Combined with gauge for all-hard \\
Step 9: Ablation & Manuscript \S3.2--3.5 & Main results section \\
Step 10: UQ & Manuscript \S3.6 & Connects to original paper's ensemble finding \\
Step 11: Interp & Manuscript \S3.7 & New scientific knowledge contribution \\
\bottomrule
\end{tabular}
\end{center}

%======================================================
\section{What to Remember and Avoid}
%======================================================

\begin{warningbox}
\textbf{Critical design decisions to remember throughout:}
\begin{enumerate}[noitemsep]
  \item \textbf{LayerNorm everywhere, not BatchNorm.} Hard constraint
    variants with small batches may expose BatchNorm instability.
    LayerNorm is batch-size independent.

  \item \textbf{Softplus reparameterization in forward pass, not stored.}
    Non-negative weights for the pressure pathway must be recomputed
    at each forward pass via \texttt{F.softplus(W\_raw)}.
    Do not store \texttt{softplus(W)} as a parameter.

  \item \textbf{Softplus or sigmoid only on the pressure pathway.}
    Swish activations are globally used in compositional encoders but
    must not be used on the pressure pathway in V4/V7, as Swish is
    non-monotone near $x \approx -1.3$. Use softplus on all
    pressure-dependent connections.

  \item \textbf{Pressure must not enter any unconstrained pathway (V4, V7).}
    If pressure $P$ is passed through any standard linear layer with
    unrestricted weights, the monotonicity guarantee is lost entirely.
    Verify at implementation time that the only pathway from $P$ to the
    output is the reparameterized non-negative pathway.

  \item \textbf{Assert $P > 0$ before every Sips gauge computation.}
    The Sips monotonicity guarantee requires physical pressure to be
    strictly positive. NaN gradients will occur at $P = 0$.

  \item \textbf{V5 is not a single-constraint isolate.} The Sips gauge
    enforces bounds and monotonicity implicitly. Acknowledge this explicitly
    in the manuscript when interpreting V5's main effect. This is not a
    weakness --- it is the architectural consequence of choosing a
    physically structured output layer.

  \item \textbf{Identical seeds and hyperparameters across variants.}
    Any difference in initialization or training hyperparameters
    confounds the comparison. Only the constraint enforcement strategy
    changes between variants.

  \item \textbf{Report negative results honestly.} If all-hard (V2)
    underperforms all-soft (V1) on interpolation accuracy, this is a
    publishable finding. The paper's value is the rigorous comparison.

  \item \textbf{Bonferroni correction is mandatory.} With 28 pairwise
    comparisons, uncorrected $p$-values at $\alpha=0.05$ produce
    approximately 1.4 false positives by chance. Corrected $\alpha=0.0018$.

  \item \textbf{Extrapolation results are the primary contribution.}
    If hard and soft constraints perform similarly on interpolation
    but hard constraints dominate on T2/T3, the abstract and conclusions
    lead with this finding.
\end{enumerate}
\end{warningbox}

%======================================================
\section{Project Report Structure}
%======================================================

\subsection{Tentative Title}

``Guaranteeing Physical Validity by Construction: A Systematic
Factorial Evaluation of Hard vs.\ Soft Thermodynamic Constraint
Enforcement in Physics-Informed Neural Networks for Sorption Studies''

\subsection{Report Structure}

\begin{enumerate}[label=\textbf{\arabic*.}]

\item \textbf{Introduction}
\begin{itemize}
  \item Physical validity in safety-critical ML: why violations matter
  \item Soft constraint limitations: tuning-dependent, violation-tolerant,
    extrapolation-fragile
  \item Hard constraint PINNs: state-of-art in PDEs, gap in sorption
  \item Three architectural strategies and their subsumption relationships
  \item Three research questions and six contributions
\end{itemize}

\item \textbf{Materials and Methods}
\begin{itemize}
  \item Dataset: 993 measurements, 114 experiments (reference original)
  \item Three-tier evaluation framework (T1/T2/T3) with rationale
  \item Soft constraint PINN baseline: architecture + $\lambda_p$ formulation
  \item Hard bound enforcement: sigmoid-bounded output activation
  \item Hard monotonicity: non-negative pressure pathway with monotone activations
  \item Hard Sips: gauge function parameterization; subsumption of bounds and mono
  \item Hard van't~Hoff: Arrhenius structural reparameterization
  \item Revised factorial ablation: 8 variants, two clean $2\times2$ interaction subspaces
  \item Unified evaluation harness: all metrics, bootstrap CIs, McNemar, ANOVA
  \item UQ: MC Dropout + temperature scaling for all variants
  \item Interpretability: SHAP per parameter head, ALE, physical validation
\end{itemize}

\item \textbf{Results and Discussion}
\begin{itemize}
  \item \textbf{3.1} $\lambda_p$ sensitivity: soft constraint performance
    ceiling with optimal tuning
  \item \textbf{3.2} Main effects: interpolation comparison (T1), 8 variants
  \item \textbf{3.3} Physical violation rates: hard constraints achieve
    zero; statistical comparison
  \item \textbf{3.4} Extrapolation analysis (T2, T3): hard constraint
    advantage under distribution shift
  \item \textbf{3.5} Constraint interaction effects: ANOVA results on
    two clean $2\times2$ subspaces; superadditive/antagonistic findings
  \item \textbf{3.6} Uncertainty calibration: epistemic fraction vs.\
    constraint hardness; calibration under extrapolation
  \item \textbf{3.7} Interpretability: gauge parameter distributions,
    $\Delta\hat{H}$ validation, SHAP per parameter head
  \item \textbf{3.8} Computational cost analysis: Pareto frontier
    (accuracy / violation rate / training time)
  \item \textbf{3.9} Negative results and limitations (honest reporting)
\end{itemize}

\item \textbf{Conclusions}

\item \textbf{Repository} [Closed initially; structured with scripts,
run-files, reports, checkpoints, figures, and metrics, with a readme.
To be considered for open access after review.]

\end{enumerate}

\end{document}
